\documentclass[10pt,a4paper]{amsart}
\usepackage{amsmath,amssymb,amsthm,graphicx}
\numberwithin{equation}{section}

\newtheorem{Theorem}{Theorem}[section]
\newtheorem{Corollary}[Theorem]{Corollary}
\newtheorem{Lemma}[Theorem]{Lemma}
\newtheorem{Proposition}[Theorem]{Proposition}
{\theoremstyle{definition}
\newtheorem{Note}[Theorem]{Note}
\newtheorem{Definition}[Theorem]{Definition}
\newtheorem{Problem}[Theorem]{Problem}
}
\usepackage[a4paper,margin=22mm]{geometry}
\title{The commutator ideal of the universal Askey-Wilson algebra equals the two-sided ideal generated by the commutator of A and B inside the subalgebra generated by A and B}
\author{Paul Terwilliger}
\date{doi:10.3842/SIGMA.2011.069; CC BY 4.0}
\begin{document}
\maketitle
\noindent\textit{Editorial note.} The counted claim is the author Lemma lem:comu at source lines 3123-3131, and its proof is the proof environment that immediately follows it at source lines 3133-3168; the earlier lemmas are retained because the counted proof invokes them by label.\par\medskip
\noindent\textit{Editorial note.} The spans at source lines 2467-2474 and 3009-3016 are retained so that the abbreviation O = <A,B> used by the counted statement is defined inside the material; the surrounding Section 10 discussion of the homomorphism from the q-Onsager algebra is not part of this task.\par\medskip
\noindent\textit{Editorial note.} eq:caspreview is a display inside the introduction and is retained only because the definition of alpha, beta, gamma refers to it; it is not used as an input to any step of the counted proof.\par\medskip
\noindent\textit{Allowed standard prerequisites (editorial).}\begin{itemize}
\item Associative algebra and ideal conventions. Algebras are associative with a unit, subalgebras share that unit, and the two-sided ideal generated by a subset is the span of all products u s v with s in the subset.. Used in lem:three, lem:ideal and lem:comu at source lines 3049-3168.
\item Commutativity criterion for a quotient algebra. If the images of a generating set commute in a quotient algebra, the quotient is commutative, so every commutator lies in the kernel.. Used at source lines 3066-3079 for the containment of the commutator subspace of O.
\item Centrality of gamma. The element gamma of def:abc is central in Delta, as recorded in the retained definition at source lines 301-330.. Used at source lines 3094-3095, 3114-3117 and 3146-3150.
\end{itemize}\medskip
Our conventions for the paper are as follows.
An algebra is meant to be associative and have a 1.
A subalgebra has the same 1 as the parent algebra.
We f\/ix a f\/ield $\mathbb F$
and  a nonzero $q \in \mathbb F$ such that $q^4\not=1$.

\begin{Definition}
\label{def:uaw}
Def\/ine an $\mathbb F$-algebra
$\Delta=\Delta_q$
 by generators and relations in the following way.
The generators are~$A$,~$B$,~$C$.
The relations assert that each of
\begin{gather}
 A+ \frac{qBC-q^{-1}CB}{q^2-q^{-2}},
%\label{eq:u1}
\qquad B+
\frac{qCA-q^{-1}AC}{q^2-q^{-2}},
%\label{eq:u2}
\qquad
C+
\frac{qAB-q^{-1}BA}{q^2-q^{-2}}
%\label{eq:u3}
\label{eq:comlist}
\end{gather}
is central in $\Delta$.
We call $\Delta$ the {\it universal Askey--Wilson algebra}.
\end{Definition}

\begin{Definition}
\label{def:abc}
\rm
For the
three central elements in
(\ref{eq:comlist}), multiply each by $q+q^{-1}$ to get~$\alpha$,~$\beta$,~$\gamma$. Thus
\begin{gather}
A+ \frac{qBC-q^{-1}CB}{q^2-q^{-2}}
=
\frac{\alpha}{q+q^{-1}},
\label{eq:u1}
\\
B+
\frac{qCA-q^{-1}AC}{q^2-q^{-2}}
=
\frac{\beta}{q+q^{-1}},
\label{eq:u2}
\\
C+
\frac{qAB-q^{-1}BA}{q^2-q^{-2}}
=
\frac{\gamma}{q+q^{-1}}.
\label{eq:u3}
\end{gather}
Note that each of $\alpha$, $\beta$, $\gamma$ is central in
$\Delta$.
(The purpose of the factor $q+q^{-1}$ is
to make the upcoming formula~(\ref{eq:caspreview})
more attractive.)
\end{Definition}

\begin{gather}
\label{eq:caspreview}
\Omega =
qABC+q^2A^2+q^{-2}B^2+q^2C^2-qA\alpha-q^{-1}B\beta -q C\gamma.
\end{gather}

\begin{Lemma}
\label{lem:abgpres}
The $\mathbb F$-algebra $\Delta $ is generated by
$A$, $B$, $\gamma$. Moreover
\begin{gather}
C  =  \frac{\gamma}{q+q^{-1}}
-\frac{qAB-q^{-1}BA}{q^2-q^{-2}},
\label{eq:getc}
\\
\alpha  =  \frac{B^2A-(q^2+q^{-2})BAB+AB^2+
(q^2-q^{-2})^2A+(q-q^{-1})^2B\gamma}{(q-q^{-1})(q^2-q^{-2})},
\label{eq:getalpha}
\\
\beta  =  \frac{A^2B-(q^2+q^{-2})ABA+BA^2+
(q^2-q^{-2})^2B+(q-q^{-1})^2A\gamma}{(q-q^{-1})(q^2-q^{-2})}.
\label{eq:getbeta}
\end{gather}
\end{Lemma}

\begin{proof}\sloppy
Line~(\ref{eq:getc})
is from (\ref{eq:u3}).
To get
(\ref{eq:getalpha}),
(\ref{eq:getbeta})
eliminate $C$ in
(\ref{eq:u1}),
(\ref{eq:u2})
using
(\ref{eq:getc}).
\end{proof}

{\samepage

\begin{Theorem}
\label{prop:altpres}
The $\mathbb F$-algebra $\Delta $ has a  presentation by generators
$A$, $B$, $\gamma$ and
relations
 \begin{gather*}
 A^3B-\lbrack 3\rbrack_q A^2BA+\lbrack 3\rbrack_q ABA^2-BA^3
= -\big(q^2-q^{-2}\big)^2(AB-BA),
\\
 B^3A-\lbrack 3\rbrack_q B^2AB+\lbrack 3\rbrack_q BAB^2-AB^3
= -\big(q^2-q^{-2}\big)^2(BA-AB),
\\
 A^2B^2-B^2A^2+\big(q^2+q^{-2}\big)(BABA-ABAB) = -\big(q-q^{-1}\big)^2(AB-BA)\gamma,
\\
  \qquad \gamma A = A \gamma,   \qquad \gamma B=B \gamma.
 \end{gather*}
\end{Theorem}

\begin{proof}
Use Lemma~\ref{lem:abgpres} to
express the def\/ining relations
for $\Delta $ in terms of $A$, $B$, $\gamma$.
\end{proof}}

\begin{Definition}
\label{def:hatvp}
Let $\phi: {\mathcal O} \to \Delta $
denote the $\mathbb F$-algebra homomorphism that sends
$X\mapsto A$ and $Y\mapsto B$. Observe that
$\phi$ is the restriction of
$ \varphi $ to~${\mathcal O}$.
\end{Definition}

\begin{Definition}
\label{def:hatu}
For notational convenience abbreviate
$\mathbb O =\langle A,B\rangle$. Thus
$\mathbb O$ is the image of~$\mathcal O$ under
the homomoprhism $\phi$
from Def\/inition~\ref{def:hatvp}.
\end{Definition}

\begin{Lemma}
\label{lem:notdir}
$\Delta = \sum_{n=0}^\infty \mathbb O \gamma^n$.
\end{Lemma}

\begin{proof}
 The
 algebra $\Delta $ is generated by
 $\mathbb O$, $\gamma$.
Moreover~$\gamma$ is central in~$\Delta $.
\end{proof}

\begin{Lemma}
\label{lem:three}
The following $(i)$--$(iii)$ hold.
\begin{enumerate}\itemsep=0pt
\item[$(i)$]
$\lbrack \mathbb O,\mathbb O\rbrack \subseteq
\mathbb O \lbrack A,B\rbrack \mathbb O$.
\item[$(ii)$]
$  \lbrack A,B\rbrack
\gamma \in
\lbrack \mathbb O,\mathbb O\rbrack $.
\item[$(iii)$]
$\mathbb O \lbrack A,B\rbrack \mathbb O \gamma \subseteq
\mathbb O \lbrack A,B\rbrack \mathbb O$.
\end{enumerate}
\end{Lemma}
\begin{proof}\sloppy
$(i)$
Abbreviate $R=
\mathbb O \lbrack A,B\rbrack \mathbb O$ and consider the
quotient algebra $\mathbb O/R$.
The ele\-ments~$A$,~$B$ generate~$\mathbb O$, and
these generators satisfy
$\lbrack A,B\rbrack \in R$.
Therefore $A+R$, $B+R$ generate
$\mathbb O/R$, and these generators commute.
This shows that $\mathbb O/R$ is commutative.
Consider the canonical map $\mathbb O \to \mathbb O/R$.
This map has kernel~$R$. The map sends
$\lbrack \mathbb O,\mathbb O\rbrack \mapsto 0 $ since~$\mathbb O/R$ is commutative. Therefore
$\lbrack \mathbb O,\mathbb O\rbrack \subseteq R$.

  $(ii)$ In the third displayed equation of
Theorem~\ref{prop:altpres},
the expression on the right is a nonzero
scalar multiple of
$\lbrack A,B \rbrack \gamma $.
The expression on the left
is equal to
$\lbrack A^2,B^2 \rbrack +(q^2+q^{-2})
\lbrack B,ABA \rbrack $
and is therefore in
$\lbrack \mathbb O,\mathbb O\rbrack $.
The result follows.

  $(iii)$ By $(i)$, $(ii)$ above and
since $\gamma$ is central.
\end{proof}

\begin{Lemma}
\label{lem:ideal}
$\mathbb O \lbrack A,B\rbrack \mathbb O$ is a $2$-sided ideal of
$\Delta $.
\end{Lemma}

\begin{proof}
Abbreviate $R =
\mathbb O \lbrack A,B\rbrack \mathbb O$.
We show $\Delta R\subseteq R$ and
$R\Delta \subseteq R$.
Recall that $\Delta $ is generated by
$\mathbb O$, $\gamma$.
By construction
 $\mathbb O R\subseteq R$ and
$R \mathbb  O\subseteq R$.
By Lemma~\ref{lem:three}$(iii)$ and
since $\gamma$ is central we have
 $\gamma R\subseteq R$ and
$R \gamma \subseteq R$.
By these comments
$\Delta R\subseteq R$ and
$R\Delta \subseteq R$.
\end{proof}

\begin{Lemma}
\label{lem:comu}
We have
\begin{gather*}
%\label{eq:ls}
\mathbb O \lbrack A,B \rbrack \mathbb O=
\Delta\lbrack  \Delta,\Delta\rbrack \Delta .
\end{gather*}
\end{Lemma}

\begin{proof}
We  have
$\mathbb O \subseteq \Delta $
and
$\lbrack A,B\rbrack \in \lbrack \Delta ,\Delta \rbrack$
so
$\mathbb O \lbrack A,B \rbrack \mathbb  O\subseteq
\Delta\lbrack  \Delta,\Delta\rbrack \Delta $.
We now show the reverse inclusion.
To this end we analyze
$\lbrack \Delta ,\Delta \rbrack$
using
Lemma~\ref{lem:notdir}.
For integers $m,n\geq 0$,
\begin{gather*}
\lbrack \mathbb  O \gamma^m,\mathbb O \gamma^n\rbrack
 =
\lbrack \mathbb  O ,\mathbb  O \rbrack \gamma^{m+n}
\overset{\text{by Lemma \ref{lem:three}({\it i})}}{\subseteq}
\mathbb  O \lbrack A,B\rbrack \mathbb O \gamma^{m+n}
%\qquad \qquad {\hbox {\rm by Lemma \ref{lem:three}(i)}}
 \overset{\text{by Lemma \ref{lem:ideal}}}{\subseteq}
\mathbb  O \lbrack A,B\rbrack \mathbb O.
%\qquad \qquad \qquad {\hbox {\rm by Lemma \ref{lem:ideal}.}}
\end{gather*}
By
this and Lemma~\ref{lem:notdir} we
obtain
$\lbrack \Delta ,\Delta \rbrack \subseteq
\mathbb  O \lbrack A,B\rbrack \mathbb O$.
Now using
Lemma~\ref{lem:ideal} we obtain
$\Delta \lbrack \Delta ,\Delta \rbrack \Delta  \subseteq
\mathbb  O \lbrack A,B\rbrack \mathbb  O$.
The result follows.
\end{proof}

\end{document}
